% ID: SEL-D0004-OB-E03
% ESTADO: validada. S005, 2026-11-07; tres pasadas del mismo autor y revisión matemática, didáctica y visual.
\textbf{Datos.}\par\nopagebreak[4]
Sea $M$ ``tener 40 años o más'' y $L$ ``ser propietario de la licencia''. La partición por edad da $P(M)=0.70$, $P(M^c)=0.30$, $P(L\mid M)=0.60$ y $P(L\mid M^c)=0.23$. El complementario de 40 o más es estrictamente menos de 40.

\textbf{a) Probabilidad total de ser propietario.}\par\nopagebreak[4]
\[
P(L)=0.70(0.60)+0.30(0.23)=0.42+0.069=\boxed{0.489}.
\]

\textbf{b) Edad sabiendo que no es propietario.}\par\nopagebreak[4]
$P(L^c)=1-0.489=0.511$ y $P(M\cap L^c)=0.70(0.40)=0.28$. Aplicamos Bayes:
\[
P(M\mid L^c)=\frac{P(M\cap L^c)}{P(L^c)}
=\frac{0.28}{0.511}=\boxed{\frac{280}{511}\approx0.547945}.
\]

\textbf{c) Propietario o menor de 40 años.}\par\nopagebreak[4]
La unión no es disjunta: los propietarios menores de 40 cuentan en ambos sucesos.
\[
P(L\cup M^c)=P(L)+P(M^c)-P(L\cap M^c)
=0.489+0.30-0.069=\boxed{0.72}.
\]
Comprobación alternativa: el complementario es ser mayor o igual de 40 y no propietario, de probabilidad 0.28; por tanto, la unión vale $1-0.28=0.72$. Las cuatro probabilidades conjuntas son $0.42,0.28,0.069,0.231$ y suman 1.
